% Section 3: Methodology
% Define the BilinearLayer math: y = (W_l x) \odot (W_r x)
% Describe datasets, models, and experimental setup

\section{Methodology}
\label{sec:methodology}

\subsection{Bilinear Layer Architecture}

The core component is the bilinear layer, which replaces standard nonlinearities with element-wise multiplication:
\begin{equation}
    \mathbf{y} = (\mathbf{W}_l \mathbf{x}) \odot (\mathbf{W}_r \mathbf{x})
    \label{eq:bilinear}
\end{equation}
where $\mathbf{W}_l, \mathbf{W}_r \in \mathbb{R}^{d_{\text{hidden}} \times d_{\text{in}}}$ are learnable weight matrices and $\odot$ denotes the Hadamard (element-wise) product.

Unlike standard MLPs with fixed nonlinearities (ReLU, GELU), bilinear layers compute \emph{input-dependent} gating: the left projection $\mathbf{W}_l \mathbf{x}$ gates the right projection $\mathbf{W}_r \mathbf{x}$.
This creates a quadratic dependence on the input that can be expressed as a third-order interaction tensor:
\begin{equation}
    y_c = \sum_{i,j} B_{cij} x_i x_j
    \label{eq:tensor}
\end{equation}
where $B_{cij} = \sum_h W_{\text{head},ch} (W_l)_{hi} (W_r)_{hj}$ and $W_{\text{head}} \in \mathbb{R}^{d_{\text{out}} \times d_{\text{hidden}}}$ is the output projection.

\subsection{Eigendecomposition for Interpretability}

For each output class $c$, we extract the $d_{\text{in}} \times d_{\text{in}}$ interaction matrix $B_c$ and symmetrize it:
\begin{equation}
    B_c^{\text{sym}} = \frac{1}{2}(B_c + B_c^\top)
\end{equation}
Symmetrization removes arbitrary asymmetry from the $W_l / W_r$ factorization while preserving the quadratic form $\mathbf{x}^\top B_c \mathbf{x}$.
The symmetric matrix admits eigendecomposition:
\begin{equation}
    B_c^{\text{sym}} = \sum_{r=1}^{d_{\text{in}}} \lambda_r^{(c)} \mathbf{v}_r^{(c)} (\mathbf{v}_r^{(c)})^\top
\end{equation}
where $\lambda_r^{(c)}$ are eigenvalues (sorted by magnitude) and $\mathbf{v}_r^{(c)} \in \mathbb{R}^{784}$ are eigenvectors.
For image inputs, eigenvectors can be reshaped to $28 \times 28$ images and visualized.

The key interpretability insight: if the eigenspectrum is \emph{low-rank} (few dominant eigenvalues), the model's decision for class $c$ depends on only a few input directions $\mathbf{v}_r^{(c)}$.
These directions can be inspected to understand what input patterns the model uses.

\subsection{Effective Rank}

We quantify spectral concentration using the ratio-based effective rank~\citep{roy2007effective}:
\begin{equation}
    \text{EffRank}(\boldsymbol{\lambda}) = \left( \frac{\|\boldsymbol{\lambda}\|_1}{\|\boldsymbol{\lambda}\|_2} \right)^2 = \frac{\left(\sum_i |\lambda_i|\right)^2}{\sum_i \lambda_i^2}
    \label{eq:effrank}
\end{equation}
This formula is scale-invariant and ranges from 1 (single dominant eigenvalue) to $d$ (uniform eigenvalues).
Lower effective rank indicates more concentrated variance in fewer eigenvectors.

\textbf{Implementation note:} An alternative entropy-based definition uses $\exp(H)$ where $H = -\sum_i p_i \log p_i$ with normalized eigenvalues $p_i = |\lambda_i|/\sum_j |\lambda_j|$.
We use the ratio-based formula (Equation~\ref{eq:effrank}) as it matches the paper's implementation and is more numerically stable.

\subsection{Language Model Analysis}
\label{sec:language_method}

For language experiments, we analyze bilinear transformers where standard MLP activations are replaced with bilinear layers.
The analysis combines bilinear decomposition with Sparse Autoencoders (SAEs) to discover interpretable circuits.

\subsubsection{Sparse Autoencoder Integration}

SAEs~\citep{bricken2023monosemanticity} decompose dense MLP activations into sparse, interpretable features:
\begin{equation}
    \mathbf{z} = \text{TopK}(W_{\text{enc}} \mathbf{h} + \mathbf{b}_{\text{enc}}), \quad \hat{\mathbf{h}} = W_{\text{dec}} \mathbf{z} + \mathbf{b}_{\text{dec}}
\end{equation}
where $\mathbf{h}$ is the MLP hidden state, $\mathbf{z}$ are sparse SAE features (only top-$k$ activations kept), and $W_{\text{dec}}$ columns are interpretable feature directions.

We use pretrained SAEs from \citet{pearce2024bilinear} with expansion ratios 4--8$\times$ and TopK sparsity ($k=30$).
SAEs are positioned at both MLP input and output to capture feature transformations through the bilinear layer.

\subsubsection{Interaction Matrix Analysis}

For each SAE output feature $o$, we compute the interaction matrix $Q^{(o)}$ between input SAE features:
\begin{equation}
    Q_{ij}^{(o)} = \sum_{\text{tokens}} z_i^{\text{in}} \cdot z_j^{\text{in}} \cdot z_o^{\text{out}}
    \label{eq:interaction}
\end{equation}
where $z_i^{\text{in}}$ are input SAE activations and $z_o^{\text{out}}$ is the output feature activation.
This weighted outer product captures which input feature \emph{pairs} most strongly activate output feature $o$.

Eigendecomposition of $Q^{(o)}$ reveals the dominant interaction patterns.
Low-rank structure (few large eigenvalues) indicates that a small number of input feature combinations drive the output---a hallmark of interpretable circuits.

\subsubsection{Low-Rank Correlation Metric}

Following the paper, we measure low-rank quality by correlating the true output activation $z_o^{\text{out}}$ with a rank-$k$ approximation:
\begin{equation}
    \hat{z}_o^{(k)} = \sum_{r=1}^{k} \lambda_r (\mathbf{z}^{\text{in}} \cdot \mathbf{v}_r)^2
\end{equation}
where $(\lambda_r, \mathbf{v}_r)$ are the top eigenvalue-eigenvector pairs of $Q^{(o)}$.
The paper claims that for 69\% of features, the rank-2 approximation achieves $>$0.75 Pearson correlation with the true activation.

\subsection{Datasets}

\subsubsection{Vision Datasets}

\textbf{MNIST}~\citep{lecun1998gradient}: 28$\times$28 grayscale handwritten digits.
Training: 60,000 images (54,000 train / 6,000 validation); Test: 10,000 images.
Images are flattened to 784-dimensional vectors and normalized to $[0, 1]$.

\textbf{Fashion-MNIST}~\citep{xiao2017fashion}: Same format as MNIST but with 10 clothing categories (T-shirt, trouser, pullover, etc.).
More challenging due to higher intra-class variability.

\textbf{EMNIST-Digits}~\citep{cohen2017emnist}: Handwritten digits from a different source distribution.
Used for cross-dataset transfer analysis in our extension.

\subsubsection{Language Models}

\textbf{Paper's primary model (ts-medium):} A 6-layer, 512-hidden-dimension bilinear transformer trained on TinyStories~\citep{eldan2023tinystories}.
Available as \texttt{tdooms/ts-medium} on HuggingFace.
\textit{Limitation:} The HuggingFace SAE repository lacks \texttt{mlp-in} SAEs for this model, preventing full Figure 8 reproduction.

\textbf{Our primary model (fw-medium):} A 16-layer, 1024-hidden-dimension bilinear transformer (335M parameters) trained on FineWeb-EDU.
Available as \texttt{tdooms/fw-medium} with complete SAE coverage (\texttt{mlp-in} and \texttt{mlp-out}).
We use layer 7 with expansion ratio 8$\times$.

\textbf{Additional model (fw-small):} A 12-layer variant (162M parameters) for cross-model comparison.
Layer 8 with expansion ratio 4$\times$.

\subsection{Experimental Setup}

\subsubsection{Vision Experiments}

We train bilinear MLPs on MNIST and Fashion-MNIST under four regularization configurations:

\begin{table}[h]
    \caption{Regularization configurations for vision experiments. Values from paper Appendix G.1.}
    \label{tab:reg_configs}
    \centering
    \begin{tabular}{lcc}
        \toprule
        Configuration & Noise Std & Weight Decay \\
        \midrule
        \texttt{none} & 0.0 & 0.0 \\
        \texttt{noise} & 0.5 & 0.0 \\
        \texttt{wd} & 0.0 & 1.0 \\
        \texttt{full} & 0.5 & 1.0 \\
        \bottomrule
    \end{tabular}
\end{table}

Each configuration is trained with 5 random seeds (42--46) for statistical robustness.
Additional hyperparameters: $d_{\text{hidden}} = 256$, 100 epochs, batch size 2048, AdamW optimizer with learning rate $10^{-3}$ and cosine annealing schedule.

\subsubsection{Language Experiments}

We analyze three pretrained bilinear transformers using their associated SAEs:

\begin{table}[h]
    \caption{Language model configurations for correlation analysis.}
    \label{tab:language_configs}
    \centering
    \begin{tabular}{lccccc}
        \toprule
        Model & Layers & Params & SAE Layer & Expansion & $k$ \\
        \midrule
        ts-medium & 6 & 29M & 4 & 4 & 30 \\
        fw-small & 12 & 162M & 8 & 4 & 30 \\
        fw-medium & 16 & 335M & 7 & 8 & 30 \\
        \bottomrule
    \end{tabular}
\end{table}

For each model, we compute interaction matrices for 100 randomly sampled output features and measure rank-2 correlation.

\subsection{Computational Setup}

Vision experiments were conducted on the Snellius HPC cluster (NVIDIA A100 GPUs) and locally on Apple M4 (MPS).
Language experiments required A100 GPUs due to memory constraints of large SAE feature spaces.
We track computational costs using \texttt{codecarbon}~\citep{codecarbon} for CO2 emission estimates.
